\label{sec:related_work}
Earlier work in QAT \citep{guassian_quant, understanding_STE} investigated different variations of the STE and demonstrate experimentally 
that the derivative of the clipped ReLU is a better approximation than the vanilla STE (derivative of identity). In practice, the former is the same as setting the gradient outside the quantization grid to zero, as per equation \eqref{eq:dl_dw_ste}, and is the most widely used formulation in QAT literature \citep{pact2018, lsq, differentiablequantization}.

Despite the success of the STE trick in practice, a lot of prior work has focused on the non-optimality of the STE as a gradient approximator. \citet{diffq} claim that the STE causes instability during training due to weights oscillating leading to \textit{biased} gradients and weights. For this reason, they propose replacing the rounding operation with additive Gaussian noise in the forward pass to simulate quantization noise. Whereas weight oscillations are indeed an important side-effect of the STE, the idea that the STE yields \textit{biased} gradients in a deterministic training setting is opaque, as we discuss in more detail in appendix \ref{app:biased_ste}. \citet{qat_abs_cosine} use a toy example to claim that solutions under STE can lead to a sub-optimal quantized solution but do not recognize that the quantized weight will oscillate around the optimum rather than converge to the wrong value. Based on their observation, they propose an absolute-cosine regularization for the full-precision weights to make them more quantization-friendly but do not compare to other competing QAT methods.

Other work focuses on alternative \textit{gradient estimators} for back-propagation. Despite seeming differences in their motivation and formulations, all these methods in on way or another propose adaptations of STE that boil down into two main categories: \textit{multiplicative} or \textit{additive}. Multiplicative methods apply a scaling to the data-gradient, whereas additive methods add an input-independent term to the gradient through regularization. We provide more details and derivations in appendix \ref{app:toy_example_gradients}. 

% Other work focuses on alternative \textit{gradient estimators} for back-propagation. Despite seeming differences in their motivation and formulations, all these methods in on way or another propose adaptations of the STE that incidentally affect the oscillatory behavior. Most of these methods boil down into two main categories: 


% Other work focuses on alternative gradient estimators for back-propagation. Despite seeming differences in their motivation and formulations, we believe that all these methods in one way or another try to address the oscillation problem by employing variations of the STE. We split these methods in two main categories:

% \vspace{-.2cm}
\paragraph{Multiplicative} \citet{EWGS} (EWGS) acknowledge that the STE generates \textit{coarse gradients} that treat all latent weights within a quantization bin equally regardless of their distance from its centre and propose scaling the gradient by this distance. However, they do not provide any evidence on why and how these coarse gradient harm optimization.  \citet{PBGS} arrive to a similar formulation but are motivated by robust quantization. They claim that clustering latent weights around bin centres makes them more robust when quantized to different bit widths. In both these methods, the gradient scales linearly with the distance of the latent weight from its quantized value.

% Motivated by the need for the latent weights to cluster around the bin centres to enable robust quantization, \citet{PBGS} arrive to a similar formulation as \citet{EWGS} (see appendix \ref{app:psg_ewgs} for more details). In both these methods, the gradient scale linearly with the rounding error. 

\citet{DSQ} (DSQ) try to alleviate the \textit{gradient error} that the STE causes during low-bit quantization but do not provide any evidence about the source or magnitude of this error. They propose using a series of hyperbolic tangent functions to simulate the rounding function in the backward pass. This leads to a non-linear scaling of the STE gradient, with the gradient increasing with the distance from the bin centre. \citet{quantization_nets} also approximate the rounding function with  a stack of sigmoid functions to address the \textit{gradient mismatch} problem. They parameterize the sigmoid with a temperature, which increases during training to better approximate the step function.

As we show in the appendix, these multiplicative methods can affect the magnitude of the oscillations but cannot stop them from happening. In figure \ref{fig:oscillation_single_weight}, we visualize the effect of EWGS and DSQ gradients on the 1D regression problem, described in section \ref{sec:oscillations}. It is evident, that neither of these methods overcome oscillations.
% add a term to the STE that scales the gradient depending of the distance of the weight from the closest grid point or \textit{rounding error}. \citet{PBGS} arrive to a similar formulation but the gradient is effectively zero when the weight is at the grid point. Both these method depend linearly on the rounding error. \citet{quantization_nets, DSQ} apply a more aggressive non-linear scaling by simulating the rounding function with the sigmoid. 
% We show in the following section that multiplicative methods can influence the amplitude of the oscillations but they cannot prevent them. \todo{add this section}
% % Whereas, the multiplicative methods can reduce amplitude or frequency of the oscillations, they cannot prevent them or dampen them. 
\vspace{-.1cm}
\paragraph{Additive} Various regularizations have been proposed in literature to make the weight distribution more quantization-friendly.
\citet{qat_tinyml} apply an $L^2$ regularization on the difference between quantized and floating-point weights, while \citet{binreg} apply bin regularization to enforce the weights to cluster at the quantization grid points. Their bin regularization method is similar to our proposed \textit{oscillation dampening} method described later in section \ref{sec:oscillations_dampening}. However, their motivation is to form sharper weight distribution around the bin centre that reduce the quantization error rather than preventing oscillations. In contrast to the multiplicative methods, these additive methods add a term to the gradient and can thus dampen oscillations by pulling the latent weights towards the grid points.

%However, strong regularization can harm the final accuracy.  

\citet{profit} recognize that the STE can cause activation instability during training of low-bit quantized MobileNets and propose progressively freezing whole layers during training. \citet{metaquant} approximate the gradient using a small neural network, whose parameters are learnt using meta-learning but is limited to binary weights.
